\section{Conclusão}
\label{sec:conclusao}

Os cinco circuitos de teste do roteiro foram simulados com o macromodelo
OP07 do LTspice e todos os parâmetros de não idealidade foram medidos:
$I_{B+}$ = \SI{0.78}{\nano\ampere}, $I_{B-}$ =
\SI{1.18}{\nano\ampere}, $I_{OS}$ =
\SI{0.40}{\nano\ampere}, $V_{OS}$ = \SI{-30.4}{\micro\volt},
$A_{vd}$ = \SI{500}{\volt\per\milli\volt}, CMRR =
\SI{126.0}{\decibel} e SR = \SI{0.260}{\volt\per\micro\second}.
Todos ficaram dentro das faixas da folha de dados e próximos dos valores
típicos do grau OP07C. As saídas dos itens 1 e 2 foram conferidas por
cálculo com os próprios parâmetros medidos, e cada valor obtido por
derivada numérica foi confirmado por um segundo estimador independente
(secante por \texttt{.meas} nos ganhos e medida clássica entre
$\pm$\SI{10}{\volt} no slew rate).

O experimento também deixou claro o papel de cada circuito de teste em
transformar um efeito muito pequeno em algo que dá para medir: os
resistores de \SI{1}{\mega\ohm} transformam nanoamperes em milivolts, a
malha aberta expõe o offset e os ganhos, a compensação do $V_{OS}$
mantém o amplificador na região linear e o degrau no seguidor isola o
slew rate. A comparação com o amplificador ideal (correntes e offset
zero, ganho, CMRR e slew rate infinitos) mostra o tamanho do desvio do
componente real em cada parâmetro.
